\Needspace{12\baselineskip}
\section{References}

\begingroup
\setlength{\parskip}{0.15em}
\titlespacing*{\subsection}{0pt}{1.35ex}{0.45ex}
\begin{multicols}{2}
\raggedcolumns

\subsection*{Scripture, Fathers, and early conciliar witnesses}

{\footnotesize
\begin{itemize}[itemsep=0.1em,topsep=0.15em]
\item Matt 16:13--19; Luke 22:31--32; John 21:15--17; Acts 1--15; Rom 12; 1 Cor 10--12; Eph 2--4; 1 Tim 3:15; 1 Pet 2 and 5. Scriptural references supply the apostolic, ecclesial, Eucharistic, and charismal grammar; no modern translation is reproduced at length.
\item St.~Clement of Rome, \emph{First Letter to the Corinthians}, especially 42--44 and 59--63; St.~Ignatius of Antioch, \emph{Smyrnaeans} 8 and \emph{Romans}, prologue; St.~Irenaeus, \emph{Against Heresies} III.3.1--3; St.~Cyprian, \emph{On the Unity of the Catholic Church}; St.~Augustine, \emph{On Baptism} and \emph{City of God}; St.~Leo the Great, letters 10, 14, and 28. These are patristic witnesses to succession, communion, Roman service, episcopacy, and doctrinal judgment, not a claim that first-millennium administration duplicated the 1870 code.
\item Councils of Ephesus (431), Chalcedon (451), Constantinople III (680--681), Constance (1414--1418), Florence, \emph{Laetentur caeli} (1439), and Trent (1545--1563), in Norman P. Tanner, ed., \emph{Decrees of the Ecumenical Councils} (Georgetown University Press, 1990). Their distinct historical and reception settings are retained.
\item \emph{Catechism of the Catholic Church}, especially nos.~74--100, 748--975, 1322--1419, and 2032--2051, \sourceurl{https://www.vatican.va/archive/ENG0015/_INDEX.HTM}{official English index}.
\end{itemize}}

\subsection*{Vatican I, its antecedents, and authoritative reception}

{\footnotesize
\begin{itemize}[itemsep=0.1em,topsep=0.15em]
\item First Vatican Council, \sourceurl{https://www.vatican.va/archive/hist_councils/i-vatican-council/index.htm}{official document index}; dogmatic constitution \emph{Dei Filius} (24 April 1870), chs.~1--4 and canons, \sourceurl{https://www.vatican.va/archive/hist_councils/i-vatican-council/documents/vat-i_const_18700424_dei-filius_la.html}{official Latin}; dogmatic constitution \emph{Pastor aeternus} (18 July 1870), chs.~1--4, \sourceurl{https://www.vatican.va/archive/hist_councils/i-vatican-council/documents/vat-i_const_18700718_pastor-aeternus_la.html}{official Latin}.
\item Bishop Vincent Gasser, \emph{The Gift of Infallibility: The Official Relatio on Infallibility of Bishop Vincent Ferrer Gasser at Vatican Council I}, trans. James T. O'Connor (Ignatius Press, 1986), especially the explanation of “personal,” “separate,” “absolute,” inquiry, and ecclesial consent; \sourceurl{https://lawcat.berkeley.edu/record/540998}{library record}.
\item German bishops, collective declaration (January--February 1875), DS 3112--3117; contemporary English diplomatic reproduction, \sourceurl{https://history.state.gov/historicaldocuments/frus1875v01/d247}{U.S. historical-document series}. Pius IX's approving reception is also recorded in later official Roman syntheses.
\item CDF, \emph{The Primacy of the Successor of Peter in the Mystery of the Church} (31 October 1998), especially nos.~3--8 and notes 21--22, \sourceurl{https://www.vatican.va/roman_curia/congregations/cfaith/documents/rc_con_cfaith_doc_19981031_primato-successore-pietro_en.html}{official English}.
\item CDF, declaration \emph{Mysterium Ecclesiae} (5 July 1973), especially nos.~2--5, \sourceurl{https://www.vatican.va/roman_curia/congregations/cfaith/documents/rc_con_cfaith_doc_19730705_mysterium-ecclesiae_en.html}{official English}; CDF, doctrinal commentary on the concluding formula of the \emph{Professio fidei} (1998), nos.~5--13, \sourceurl{https://www.vatican.va/roman_curia/congregations/cfaith/documents/rc_con_cfaith_doc_1998_professio-fidei_en.html}{official English}.
\item John Paul II, \emph{Ut unum sint} (25 May 1995), nos.~88--96; Dicastery for Promoting Christian Unity, \emph{The Bishop of Rome} (2024), status statement, synthesis, and plenary proposal, \sourceurl{https://www.christianunity.va/content/unitacristiani/en/documenti/altri-testi/the-bishop-of-rome.html}{official study page}. The latter is used as an official ecumenical study, not as a new definition.
\item Pius IX, \emph{Ineffabilis Deus} (8 December 1854); \emph{Quanta cura} and the \emph{Syllabus} (8 December 1864); \emph{Aeterni Patris} (29 June 1868); \emph{Postquam Dei munere} (20 October 1870). Pius XII, \emph{Munificentissimus Deus} (1 November 1950), \sourceurl{https://www.vatican.va/content/pius-xii/en/apost_constitutions/documents/hf_p-xii_apc_19501101_munificentissimus-deus.html}{official English}.
\item John Henry Newman, \emph{A Letter Addressed to the Duke of Norfolk} (1875), especially sections 5--7, \sourceurl{https://www.newmanreader.org/works/anglicans/volume2/gladstone/index.html}{Newman Reader witness}; used as major theological reception, not as magisterial law.
\end{itemize}}

\subsection*{Vatican II and the authority of its teaching}

{\footnotesize
\begin{itemize}[itemsep=0.1em,topsep=0.15em]
\item Second Vatican Council, \sourceurl{https://www.vatican.va/archive/hist_councils/ii_vatican_council/index.htm}{official index of all sixteen promulgated documents}. The article uses the exact documents, dates, genres, notes, and paragraph numbers rather than a free-standing “spirit.”
\item \emph{Lumen gentium} (21 November 1964), especially nos.~8, 12, 16, 18--27, 32, and 39--42, together with the Preliminary Explanatory Note and the General Secretary's \emph{Notificationes} of 16 November 1964, which transcribes the Theological Commission's declaration of 6 March 1964; \sourceurl{https://www.vatican.va/archive/hist_councils/ii_vatican_council/documents/vat-ii_const_19641121_lumen-gentium_en.html}{official English}.
\item \emph{Unitatis redintegratio} (21 November 1964), especially nos.~3--4 on imperfect communion, ecclesial elements beyond visible Catholic boundaries, and Catholic participation in ecumenism; \sourceurl{https://www.vatican.va/archive/hist_councils/ii_vatican_council/documents/vat-ii_decree_19641121_unitatis-redintegratio_en.html}{official English}.
\item \emph{Dei Verbum} (18 November 1965), especially nos.~2--12 and 21--26; \sourceurl{https://www.vatican.va/archive/hist_councils/ii_vatican_council/documents/vat-ii_const_19651118_dei-verbum_en.html}{official English}.
\item \emph{Dignitatis humanae} (7 December 1965), nos.~1--3;
\emph{Nostra aetate} (28 October 1965), nos.~2 and 4--5,
\sourceurl{https://www.vatican.va/archive/hist_councils/ii_vatican_council/documents/vat-ii_decl_19651028_nostra-aetate_en.html}{official English};
\sourceurl{https://www.vatican.va/archive/hist_councils/ii_vatican_council/documents/vat-ii_decree_19651207_ad-gentes_en.html}{\emph{Ad gentes} 2, 7} (7 December 1965); and \emph{Gaudium et spes}
(7 December 1965), nos.~17--22 (atheism, 19--21).
\item John XXIII, \emph{Gaudet Mater Ecclesia} (11 October 1962), \sourceurl{https://www.vatican.va/content/john-xxiii/la/speeches/1962/documents/hf_j-xxiii_spe_19621011_opening-council.html}{official Latin}; announcement of the council (25 January 1959), \sourceurl{https://www.vatican.va/content/john-xxiii/it/speeches/1959/documents/hf_j-xxiii_spe_19590125_annuncio.html}{official Italian}.
\item Paul VI, general audience (12 January 1966), on the council's pastoral character, avoidance of new extraordinary definitions, and authority of the supreme ordinary magisterium, \sourceurl{https://www.vatican.va/content/paul-vi/it/audiences/1966/documents/hf_p-vi_aud_19660112.html}{official Italian}; closing address (7 December 1965), \sourceurl{https://www.vatican.va/content/paul-vi/en/speeches/1965/documents/hf_p-vi_spe_19651207_epilogo-concilio.html}{official English}; apostolic brief \emph{In Spiritu Sancto} (8 December 1965), \sourceurl{https://www.vatican.va/content/paul-vi/en/apost_letters/documents/hf_p-vi_apl_19651208_in-spiritu-sancto.html}{official English}.
\item CDF, \emph{Responses to Some Questions Regarding Certain Aspects of the Doctrine on the Church} (29 June 2007), \sourceurl{https://www.vatican.va/roman_curia/congregations/cfaith/documents/rc_con_cfaith_doc_20070629_commento-responsa_en.html}{official English}; CDF, \emph{Dominus Iesus} (6 August 2000), nos.~13--17.
\item Benedict XVI, address to the Roman Curia (22 December 2005), \sourceurl{https://www.vatican.va/content/benedict-xvi/en/speeches/2005/december/documents/hf_ben_xvi_spe_20051222_roman-curia.html}{official English}, on the hermeneutic of reform and the levels at which continuity and discontinuity occur.
\end{itemize}}

\subsection*{Communism, Freemasonry, and twentieth-century renewal}

{\footnotesize
\begin{itemize}[itemsep=0.1em,topsep=0.15em]
\item Leo XIII, \emph{Humanum genus} (20 April 1884), especially nos.~9--12 and 16--29, \sourceurl{https://www.vatican.va/content/leo-xiii/en/encyclicals/documents/hf_l-xiii_enc_18840420_humanum-genus.html}{official English}.
\item Pius XI, \emph{Divini Redemptoris} (19 March 1937), \sourceurl{https://www.vatican.va/content/pius-xi/en/encyclicals/documents/hf_p-xi_enc_19370319_divini-redemptoris.html}{official English}; Holy Office, decree against Communist doctrine (1 July 1949), AAS 41 (1949), 334.
\item Leo XIII, \emph{Aeterni Patris} (4 August 1879); Pius X, \emph{Lamentabili sane} and \emph{Pascendi dominici gregis} (1907); anti-Modernist oath, \emph{Sacrorum antistitum} (1 September 1910); Pius XII, \emph{Divino afflante Spiritu} (1943) and \emph{Humani generis} (1950). These official acts govern the article's account of Roman theological policy.
\item Pius XII, \emph{Mystici Corporis Christi} (1943); Holy Office, instruction \emph{Ecclesia catholica} on the ecumenical movement (20 December 1949), AAS 42 (1950), 142--147.
\end{itemize}}

\subsection*{Liturgical reform and the Roman Missal}

{\footnotesize
\begin{itemize}[itemsep=0.1em,topsep=0.15em]
\item Pius XII, \emph{Mediator Dei} (20 November 1947), especially nos.~20--23, 59--64, and 105--109, \sourceurl{https://www.vatican.va/content/pius-xii/en/encyclicals/documents/hf_p-xii_enc_20111947_mediator-dei.html}{official English}; Sacred Congregation of Rites, \emph{Dominicae resurrectionis vigiliam} (1951) and \emph{Maxima redemptionis nostrae mysteria} (16 November 1955), official AAS texts.
\item John XXIII, \emph{Rubricarum instructum} (25 July 1960), and Sacred Congregation of Rites decree declaring the Vatican \emph{Missale Romanum} editio typica (23 June 1962). The decree's date and text were checked in the 1962 typical edition; \sourceurl{https://media.churchmusicassociation.org/pdf/missale62.pdf}{facsimile witness}.
\item Vatican II, \emph{Sacrosanctum Concilium} (4 December 1963), especially nos.~21--23, 34, 36--40, 47--58, 112, and 116, \sourceurl{https://www.vatican.va/archive/hist_councils/ii_vatican_council/documents/vat-ii_const_19631204_sacrosanctum-concilium_en.html}{official English}.
\item Paul VI, \emph{Sacram liturgiam} (25 January 1964), \sourceurl{https://www.vatican.va/content/paul-vi/la/motu_proprio/documents/hf_p-vi_motu-proprio_19640125_sacram-liturgiam.html}{official Latin}; address to the Consilium (29 October 1964), \sourceurl{https://www.vatican.va/content/paul-vi/la/speeches/1964/documents/hf_p-vi_spe_19641029_consilium.html}{official Latin}.
\item Sacred Congregation of Rites and Consilium, \emph{Inter Oecumenici} (26 September 1964), AAS 56 (1964), 877--900; \emph{Musicam sacram} (5 March 1967), \sourceurl{https://www.vatican.va/archive/hist_councils/ii_vatican_council/documents/vat-ii_instr_19670305_musicam-sacram_en.html}{official English}; \emph{Tres abhinc annos} (4 May 1967), AAS 59 (1967), 442--448.
\item Paul VI, \emph{Mysterium fidei} (3 September 1965), \sourceurl{https://www.vatican.va/content/paul-vi/en/encyclicals/documents/hf_p-vi_enc_03091965_mysterium.html}{official English}; \emph{Mysterii paschalis} (14 February 1969), \sourceurl{https://www.vatican.va/content/paul-vi/en/motu_proprio/documents/hf_p-vi_motu-proprio_19690214_mysterii-paschalis.html}{official English}.
\item Paul VI, apostolic constitution \emph{Missale Romanum} (3 April 1969), \sourceurl{https://www.vatican.va/content/paul-vi/en/apost_constitutions/documents/hf_p-vi_apc_19690403_missale-romanum.html}{official English}; general audiences of \sourceurl{https://www.vatican.va/content/paul-vi/it/audiences/1969/documents/hf_p-vi_aud_19691119.html}{19 November} and \sourceurl{https://www.vatican.va/content/paul-vi/it/audiences/1969/documents/hf_p-vi_aud_19691126.html}{26 November 1969}, official Italian.
\item Sacred Congregation for Divine Worship, \emph{Memoriale Domini} (29 May 1969), AAS 61 (1969), 541--547, \sourceurl{https://www.vatican.va/archive/aas/documents/AAS-61-1969-ocr.pdf}{official AAS volume}. It governs the article's narrow historical claim about retained usage and a territorially regulated contrary practice.
\item \emph{Missale Romanum}, editiones typicae 1970, 1975, and 2002; General Instruction of the Roman Missal, third typical edition, \sourceurl{https://www.vatican.va/content/dam/wss/roman_curia/congregations/ccdds/documents/rc_con_ccdds_doc_20030317_ordinamento-messale_en.html}{official English}; Holy See press conference for the third typical edition (22 March 2002), \sourceurl{https://press.vatican.va/content/salastampa/it/bollettino/pubblico/2002/03/22/0150/00449.html}{official Italian}.
\item Congregation for Divine Worship, \emph{Liturgiam authenticam} (28 March 2001), \sourceurl{https://press.vatican.va/roman_curia/congregations/ccdds/documents/rc_con_ccdds_doc_20010507_liturgiam-authenticam_en.html}{official English}; used for the typical-edition and translation distinction, not as a complete history of each territory.
\end{itemize}}

\subsection*{Religious life, practice, and institutional data}

{\footnotesize
\begin{itemize}[itemsep=0.1em,topsep=0.15em]
\item Vatican II, \emph{Perfectae caritatis} (28 October 1965), especially nos.~2--6, 12--15, and 17, \sourceurl{https://www.vatican.va/archive/hist_councils/ii_vatican_council/documents/vat-ii_decree_19651028_perfectae-caritatis_en.html}{official English}; Paul VI, \emph{Ecclesiae Sanctae} (6 August 1966), section II, parts I--II, \sourceurl{https://www.vatican.va/content/paul-vi/en/motu_proprio/documents/hf_p-vi_motu-proprio_19660806_ecclesiae-sanctae.html}{official English}.
\item Congregation for Institutes of Consecrated Life, \emph{Report on the Apostolic Visitation of Institutes of Women Religious in the United States} (16 December 2014), \sourceurl{https://www.usccb.org/beliefs-and-teachings/vocations/consecrated-life/apostolic-visitation-final-report}{USCCB-hosted official report page}; used for the exceptional postwar-entry peak and visitation findings.
\item CARA, \emph{Frequently Requested Church Statistics}, current workbook, United States series through 2025 and worldwide series through 2023, \sourceurl{https://cara.georgetown.edu/faqs}{source page}. United States counts derive principally from the \emph{Official Catholic Directory}; worldwide counts derive from Vatican statistical yearbooks. Exact rows, denominators, and calculations are recorded in the repository research inventory.
\item Holy See Press Office / Fides, \emph{Catholic Church Statistics 2025} (2023 data), \sourceurl{https://press.vatican.va/content/dam/salastampa/it/fuori-bollettino/pdf/EN\%20-\%20Catholic\%20Church\%20Statistics\%202025.pdf}{official PDF}, used as an independent official check on current global scale and geography.
\item National Religious Vocation Conference / CARA, \emph{Recent Vocations to Religious Life in the United States} (2020), \sourceurl{https://nrvc.net/509/publication/9180/article/21258-2020-study-overview-and-links-to-study-resources}{study portal}, on characteristics reported by recent entrants and institutes.
\item John Jay College Research Team, \emph{The Causes and Context of Sexual Abuse of Minors by Catholic Priests in the United States, 1950--2010} (2011), \sourceurl{https://www.usccb.org/sites/default/files/issues-and-action/child-and-youth-protection/upload/The-Causes-and-Context-of-Sexual-Abuse-of-Minors-by-Catholic-Priests-in-the-United-States-1950-2010.pdf}{report PDF}. Its data and multi-cause conclusion are used with the report's stated limitations.
\end{itemize}}

\subsection*{Belief and identity surveys}

{\footnotesize
\begin{itemize}[itemsep=0.1em,topsep=0.15em]
\item Pew Research Center, \emph{U.S. Protestants Are Not Defined by Reformation-Era Controversies} (31 August 2017), Catholic subsample and exact forced choices, \sourceurl{https://www.pewresearch.org/religion/2017/08/31/u-s-protestants-are-not-defined-by-reformation-era-controversies-500-years-later/}{report}.
\item Pew Research Center, \emph{Just one-third of U.S. Catholics agree with their church that Eucharist is body, blood of Christ} (5 August 2019), question wording, knowledge cross-tabulation, and attendance, \sourceurl{https://www.pewresearch.org/short-reads/2019/08/05/transubstantiation-eucharist-u-s-catholics/}{report}.
\item CARA / NORC, \emph{Eucharist: National Survey of Catholic Adults} (fielded 11 July--2 August 2022; published September 2023), sample, open and closed items, coding, and margin of error, \sourceurl{https://cara.georgetown.edu/s/EucharistPollSeptember23.pdf}{report PDF}.
\item Pew Research Center, \emph{Essentials of Catholic Identity} (16 June 2025), \sourceurl{https://www.pewresearch.org/religion/2025/06/16/essentials-of-catholic-identity/}{report}; Pew, \emph{Most U.S. Catholics Say They Want the Church to Be More Inclusive} (30 April 2025), \sourceurl{https://www.pewresearch.org/religion/2025/04/30/most-us-catholics-say-they-want-the-church-to-be-more-inclusive/}{report}.
\item Pew Research Center, seven-country Catholic survey (26 September 2024), \sourceurl{https://www.pewresearch.org/religion/2024/09/26/many-catholics-in-the-us-and-latin-america-want-the-church-to-allow-birth-control-and-to-let-women-become-priests/}{report}. Policy questions are never relabeled as direct dogmatic-assent measures.
\end{itemize}}

\subsection*{SSPX, penalties, and older Roman books}

{\footnotesize
\begin{itemize}[itemsep=0.1em,topsep=0.15em]
\item 1983 \emph{Code of Canon Law}, can.~751; current cann.~1311--1363, 1364, and 1387; Francis, \emph{Pascite gregem Dei} (23 May 2021), promulgating revised Book VI effective 8 December 2021, \sourceurl{https://www.vatican.va/content/francesco/en/apost_constitutions/documents/papa-francesco_costituzione-ap_20210523_pascite-gregem-dei.html}{official English}. Unauthorized episcopal consecration stood at can.~1382 before revision and at can.~1387 afterward.
\item John Paul II, \emph{Ecclesia Dei adflicta} (2 July 1988), nos.~3--6, \sourceurl{https://www.vatican.va/content/john-paul-ii/en/motu_proprio/documents/hf_jp-ii_motu-proprio_02071988_ecclesia-dei.html}{official English}; Pontifical Council for Legislative Texts, explanatory note (24 August 1996), nos.~5--7, \sourceurl{https://www.vatican.va/roman_curia/pontifical_councils/intrptxt/documents/rc_pc_intrptxt_doc_19960824_vescovo-lefebvre_it.html}{official Italian}.
\item Congregation for Bishops, decree remitting the 1988 censures (21 January 2009), \sourceurl{https://www.vatican.va/roman_curia/congregations/cbishops/documents/rc_con_cbishops_doc_20090121_remissione-scomunica_en.html}{official English}; Benedict XVI, explanatory letter (10 March 2009), \sourceurl{https://www.vatican.va/content/benedict-xvi/en/letters/2009/documents/hf_ben-xvi_let_20090310_remissione-scomunica.pdf}{official PDF}; \emph{Ecclesiae unitatem} (2 July 2009).
\item Francis, \emph{Misericordia et misera} (20 November 2016), no.~12, \sourceurl{https://www.vatican.va/content/francesco/en/apost_letters/documents/papa-francesco-lettera-ap_20161120_misericordia-et-misera.pdf}{official PDF}; Pontifical Commission \emph{Ecclesia Dei}, letter on SSPX marriages (27 March / published 4 April 2017), \sourceurl{https://press.vatican.va/content/salastampa/it/bollettino/pubblico/2017/04/04/0218/00485.html}{official bulletin}.
\item DDF, statement warning that the announced ordinations would constitute a schismatic act (13 May 2026), \sourceurl{https://www.vatican.va/roman_curia/congregations/cfaith/documents/rc_ddf_doc_20260513_dichiarazione-fernandez-fsspx_en.html}{official English}; Leo XIV, letter to the SSPX superior general (29 June 2026), \sourceurl{https://www.vatican.va/content/leo-xiv/en/letters/2026/documents/20260629-lettera-fraternita-sanpiox.html}{official English}.
\item DDF, signed decree (2 July 2026), \sourceurl{https://www.vatican.va/roman_curia/congregations/cfaith/documents/rc_ddf_doc_20260702_decreto-scomunica-fsspx_it.html}{official Italian decree}, controlling the five reserved canon 1387 cases and Fellay's separate canon 1364 \S1 case; DDF, \sourceurl{https://www.vatican.va/roman_curia/congregations/cfaith/documents/rc_ddf_doc_20260702_nota-esplicativa-fsspx_it.html}{official Italian explanatory note}, controlling the current minister, lay-adherence, and sacramental consequences. \sourceurl{https://www.vaticannews.va/en/church/news/2026-07/lefebvrists-consecrate-four-new-bishops-without-a-papal-mandate.html}{Vatican News's 1 July event report} is used only to establish and publicly report the completed event, not to allocate penalties.
\item Congregation for Divine Worship, \emph{Quattuor abhinc annos} (3 October 1984), AAS 76 (1984), 1088--1089; Benedict XVI, \emph{Summorum Pontificum} (7 July 2007), \sourceurl{https://www.vatican.va/content/benedict-xvi/en/motu_proprio/documents/hf_ben-xvi_motu-proprio_20070707_summorum-pontificum.html}{official English}; Pontifical Commission \emph{Ecclesia Dei}, \emph{Universae Ecclesiae} (30 April 2011).
\item Francis, \emph{Traditionis custodes} (16 July 2021), \sourceurl{https://press.vatican.va/content/salastampa/it/bollettino/pubblico/2021/07/16/0469/01014.html}{official text}; CDW, Responsa (4 December / published 18 December 2021), \sourceurl{https://www.vatican.va/roman_curia/congregations/ccdds/documents/rc_con_ccdds_doc_20211204_responsa-ad-dubia-tradizionis-custodes_en.html}{official English}; rescript (20 February 2023), \sourceurl{https://www.vatican.va/roman_curia/congregations/ccdds/documents/rc_con_ccdds_doc_20230220_rescriptum-traditioniscustodes_en.html}{official English}; particular FSSP decree (11 February 2022), \sourceurl{https://www.fssp.pt/documentos/decretopapafrancisco}{institute-hosted facsimile}.
\end{itemize}}

\subsection*{Historical, liturgical, and sociological studies}

{\footnotesize
\begin{itemize}[itemsep=0.1em,topsep=0.15em]
\item John W. O'Malley, \emph{Vatican I: The Council and the Making of the Ultramontane Church} (Harvard University Press, 2018), \sourceurl{https://books.google.com/books/about/Vatican_I.html?id=eSRgDwAAQBAJ}{bibliographic preview}; Klaus Schatz, \emph{Papal Primacy: From Its Origins to the Present} (Liturgical Press, 1996); Hermann J. Pottmeyer, \emph{Towards a Papacy in Communion} (Crossroad, 1998).
\item John W. O'Malley, \emph{What Happened at Vatican II} (Harvard University Press, 2008); Joseph Ratzinger, \emph{Theological Highlights of Vatican II} (Paulist Press, 1966); Yves Congar, \emph{My Journal of the Council} (Liturgical Press, 2012). Participant theological records are checked against promulgated texts.
\item Josef A. Jungmann, \emph{The Mass of the Roman Rite}; Annibale Bugnini, \emph{The Reform of the Liturgy 1948--1975} (Liturgical Press, 1990), \sourceurl{https://books.google.com/books/about/The_Reform_of_the_Liturgy_1948_1975.html?id=FUfZAAAAMAAJ}{bibliographic record}; Alcuin Reid, \emph{The Organic Development of the Liturgy} (Ignatius Press, 2005).
\item Paul F. Bradshaw, Maxwell E. Johnson, and L. Edward Phillips, \emph{The Apostolic Tradition: A Commentary} (Fortress Press, 2002), used to bound claims about the disputed textual history behind Eucharistic Prayer II; Lauren Pristas, \emph{Collects of the Roman Missals} (Bloomsbury/T\&T Clark, 2013), for comparative oration method.
\item Stephen Bullivant, \emph{Mass Exodus: Catholic Disaffiliation in Britain and America since Vatican II} (Oxford University Press, 2019), \sourceurl{https://academic.oup.com/book/35175}{publisher record}; Patricia Wittberg, \emph{The Rise and Fall of Catholic Religious Orders} (SUNY Press, 1994), \sourceurl{https://utpdistribution.com/9780791422304/the-rise-and-fall-of-catholic-religious-orders/}{publisher/distributor record}.
\item Rodney Stark and Roger Finke, \emph{Acts of Faith: Explaining the Human Side of Religion} (University of California Press, 2000). Its institutional incentive model is used as one sociological lens and not as a theology of supernatural vocation.
\end{itemize}}

\end{multicols}
\endgroup
